% Options for packages loaded elsewhere
% Options for packages loaded elsewhere
\PassOptionsToPackage{unicode}{hyperref}
\PassOptionsToPackage{hyphens}{url}
%
\documentclass[
  14pt,
  ignorenonframetext,
  aspectratio=169,
]{beamer}
\newif\ifbibliography
\usepackage{pgfpages}
\setbeamertemplate{caption}[numbered]
\setbeamertemplate{caption label separator}{: }
\setbeamercolor{caption name}{fg=normal text.fg}
\beamertemplatenavigationsymbolshorizontal
% remove section numbering
\setbeamertemplate{part page}{
  \centering
  \begin{beamercolorbox}[sep=16pt,center]{part title}
    \usebeamerfont{part title}\insertpart\par
  \end{beamercolorbox}
}
\setbeamertemplate{section page}{
  \centering
  \begin{beamercolorbox}[sep=12pt,center]{section title}
    \usebeamerfont{section title}\insertsection\par
  \end{beamercolorbox}
}
\setbeamertemplate{subsection page}{
  \centering
  \begin{beamercolorbox}[sep=8pt,center]{subsection title}
    \usebeamerfont{subsection title}\insertsubsection\par
  \end{beamercolorbox}
}
% Prevent slide breaks in the middle of a paragraph
\widowpenalties 1 10000
\raggedbottom
\AtBeginPart{
  \frame{\partpage}
}
\AtBeginSection{
  \ifbibliography
  \else
    \frame{\sectionpage}
  \fi
}
\AtBeginSubsection{
  \frame{\subsectionpage}
}
\usepackage{iftex}
\ifPDFTeX
  \usepackage[T1]{fontenc}
  \usepackage[utf8]{inputenc}
  \usepackage{textcomp} % provide euro and other symbols
\else % if luatex or xetex
  \usepackage{unicode-math} % this also loads fontspec
  \defaultfontfeatures{Scale=MatchLowercase}
  \defaultfontfeatures[\rmfamily]{Ligatures=TeX,Scale=1}
\fi
\usepackage{lmodern}

\usetheme[]{metropolis}
\usecolortheme[]{default}
\ifPDFTeX\else
  % xetex/luatex font selection
\fi
% Use upquote if available, for straight quotes in verbatim environments
\IfFileExists{upquote.sty}{\usepackage{upquote}}{}
\IfFileExists{microtype.sty}{% use microtype if available
  \usepackage[]{microtype}
  \UseMicrotypeSet[protrusion]{basicmath} % disable protrusion for tt fonts
}{}
\makeatletter
\@ifundefined{KOMAClassName}{% if non-KOMA class
  \IfFileExists{parskip.sty}{%
    \usepackage{parskip}
  }{% else
    \setlength{\parindent}{0pt}
    \setlength{\parskip}{6pt plus 2pt minus 1pt}}
}{% if KOMA class
  \KOMAoptions{parskip=half}}
\makeatother


\usepackage{longtable,booktabs,array}
\newcounter{none} % for unnumbered tables
\usepackage{calc} % for calculating minipage widths
\usepackage{caption}
% Make caption package work with longtable
\makeatletter
\def\fnum@table{\tablename~\thetable}
\makeatother
\usepackage{graphicx}
\makeatletter
\newsavebox\pandoc@box
\newcommand*\pandocbounded[1]{% scales image to fit in text height/width
  \sbox\pandoc@box{#1}%
  \Gscale@div\@tempa{\textheight}{\dimexpr\ht\pandoc@box+\dp\pandoc@box\relax}%
  \Gscale@div\@tempb{\linewidth}{\wd\pandoc@box}%
  \ifdim\@tempb\p@<\@tempa\p@\let\@tempa\@tempb\fi% select the smaller of both
  \ifdim\@tempa\p@<\p@\scalebox{\@tempa}{\usebox\pandoc@box}%
  \else\usebox{\pandoc@box}%
  \fi%
}
% Set default figure placement to htbp
\def\fps@figure{htbp}
\makeatother





\setlength{\emergencystretch}{3em} % prevent overfull lines

\providecommand{\tightlist}{%
  \setlength{\itemsep}{0pt}\setlength{\parskip}{0pt}}



 


\setbeamertemplate{footline}[frame number]
\setbeamertemplate{section in toc}[sections numbered]
\setbeamertemplate{frametitle continuation}{}
\makeatletter
\@ifpackageloaded{caption}{}{\usepackage{caption}}
\AtBeginDocument{%
\ifdefined\contentsname
  \renewcommand*\contentsname{Table of contents}
\else
  \newcommand\contentsname{Table of contents}
\fi
\ifdefined\listfigurename
  \renewcommand*\listfigurename{List of Figures}
\else
  \newcommand\listfigurename{List of Figures}
\fi
\ifdefined\listtablename
  \renewcommand*\listtablename{List of Tables}
\else
  \newcommand\listtablename{List of Tables}
\fi
\ifdefined\figurename
  \renewcommand*\figurename{Figure}
\else
  \newcommand\figurename{Figure}
\fi
\ifdefined\tablename
  \renewcommand*\tablename{Table}
\else
  \newcommand\tablename{Table}
\fi
}
\@ifpackageloaded{float}{}{\usepackage{float}}
\floatstyle{ruled}
\@ifundefined{c@chapter}{\newfloat{codelisting}{h}{lop}}{\newfloat{codelisting}{h}{lop}[chapter]}
\floatname{codelisting}{Listing}
\newcommand*\listoflistings{\listof{codelisting}{List of Listings}}
\makeatother
\makeatletter
\makeatother
\makeatletter
\@ifpackageloaded{caption}{}{\usepackage{caption}}
\@ifpackageloaded{subcaption}{}{\usepackage{subcaption}}
\makeatother

\usepackage{bookmark}
\IfFileExists{xurl.sty}{\usepackage{xurl}}{} % add URL line breaks if available
\urlstyle{same}
\hypersetup{
  pdftitle={Clinical Trials Workshop},
  pdfauthor={Prof.~A. Nabhan},
  pdfsubject={Research},
  hidelinks,
  pdfcreator={LaTeX via pandoc}}


\title{\texorpdfstring{Clinical Trials}{Clinical Trials Workshop}}
\subtitle{\texorpdfstring{A Hands-on Workshop}{A Hands-on Workshop}}
\author{Prof.~Ashraf Nabhan}
\date{}

\begin{document}
\frame{\titlepage}

\renewcommand*\contentsname{What we will do today}
\begin{frame}[allowframebreaks]
  \frametitle{What we will do today}
  \setcounter{tocdepth}{3}
  \tableofcontents
\end{frame}
\setcounter{tocdepth}{3}
\tableofcontents
}

\section{Clinical trial Methods: A
Recap}\label{clinical-trial-methods-a-recap}

\begin{frame}{The essentials}
\protect\phantomsection\label{the-essentials}
\begin{itemize}
\item
  Type of trial: parallel, multi-arm, cross-over, etc
\item
  Unit of allocation: individual, cluster
\item
  Framework: superiority, non-inferiority
\item
  Allocation: Non-random, Quasi-random, Random
\item
  Controlled or not
\item
  Phase: Only for drug trials
\end{itemize}
\end{frame}

\section{Sample size calculation}\label{sample-size-calculation}

\begin{frame}{Preample}
\protect\phantomsection\label{preample}
\begin{itemize}
\item
  Estimated number of participants needed to achieve study objectives
\item
  Describe how it was determined
\item
  Explicitly indicate clinical and statistical assumptions supporting
  calculations.
\end{itemize}
\end{frame}

\begin{frame}{Ingredients}
\protect\phantomsection\label{ingredients}
\begin{itemize}
\item
  treat: the expected proportion for the treatment group
\item
  control: the expected proportion for the control group
\item
  power: the required study power
\item
  r: the number in the treatment group divided by the number in the
  control group.
\item
  sided.test: Use a two-sided test if you wish to evaluate whether the
  outcome proportion in the treatment group is greater than or less than
  the outcome proportion in the control group.
\item
  conf.level: defining the level of confidence in the computed result.
\end{itemize}
\end{frame}

\begin{frame}{Practical exercise}
\protect\phantomsection\label{practical-exercise}
\begin{itemize}
\tightlist
\item
  Clinical question: Among ICU patients with septic shock and AKI
  without urgent need for dialysis, does an early, compared to delayed,
  initiation of renal replacement therapy reduce all-cause mortality at
  90 days?
\end{itemize}
\end{frame}

\begin{frame}{Ingredients}
\protect\phantomsection\label{ingredients-1}
\begin{itemize}
\item
  treat = .4
\item
  control = .5
\item
  power = .8
\item
  r = 1
\item
  sided.test = 2
\item
  conf.level = .95
\end{itemize}
\end{frame}

\begin{frame}{Calculation}
\protect\phantomsection\label{calculation}
\begin{itemize}
\item
  The estimated sample size is 192
\item
  Treatment arm: 96
\item
  Control arm: 96
\end{itemize}
\end{frame}

\section{Random allocation}\label{random-allocation}

\begin{frame}{Simple random allocation}
\protect\phantomsection\label{simple-random-allocation}
\begin{itemize}
\item
  Equal probability
\item
  Unpredictability
\item
  No guarantee of balance
\end{itemize}
\end{frame}

\begin{frame}{Random allocation rule}
\protect\phantomsection\label{random-allocation-rule}
\begin{itemize}
\item
  Forced equal allocation
\item
  Vulnerability to selection bias late in trial
\end{itemize}
\end{frame}

\begin{frame}{Block randomization}
\protect\phantomsection\label{block-randomization}
\begin{itemize}
\item
  Balanced allocation within blocks
\item
  Protection against temporal trends
\item
  Predictability risk with fixed block sizes
\end{itemize}
\end{frame}

\section{Clinical trial registration}\label{clinical-trial-registration}

\begin{frame}{Open access primary registries}
\protect\phantomsection\label{open-access-primary-registries}
\begin{itemize}
\item
  Example

  \begin{itemize}
  \item
    PACTR
  \item
    Clinicaltrials.gov
  \end{itemize}
\item
  Mandatory
\item
  Use SPIRIT for the protocol

  \begin{itemize}
  \tightlist
  \item
    Perfectly aligns with trial registries
  \end{itemize}
\end{itemize}
\end{frame}

\section{Learning from mistakes}\label{learning-from-mistakes}

\begin{frame}{Comments: RCT}
\protect\phantomsection\label{comments-rct}
\begin{enumerate}
\item
  Selection bias
\item
  Serious contradiction in retinopathy inclusion criterion
\item
  Only one HbA1c measurement at the end of follow up
\item
  Intervention was not standardized
\item
  Composite outcome problem
\item
  Follow-up Duration
\end{enumerate}
\end{frame}

\section{Discussion and ideas}\label{discussion-and-ideas}




\end{document}
